%% toc: P3 38 38.1,38.2,38.3
\section{Proofs for the non-expressibility of connectedness and compactness}
\label{sec:annex-p03-w03-p03-051-nondefinability}

Throughout, the assumptions of \cref{def:p03-w03-p03-051-assumptions} are in
force. A two-sorted structure \(S=(X',O',E',(Q'_p)_p)\) for \(L_2\) has
nonempty sorts \(X'\) and \(O'\), a relation \(E'\subseteq X'\times O'\)
interpreting \(\varepsilon\), and subsets \(Q'_p\subseteq X'\)
interpreting \(P_p\); each equality symbol is identity on its own sort.
A topological or basoid model is the special case in which \(O'\) is a
family of subsets of \(X'\) and \(E'\) is membership.

\begin{annexitem}{annex:p03-w03-p03-051-two-sorted-ls}{thm:p03-w03-p03-051-two-sorted-ls}{Encoding and decoding two sorts}
\begin{proof}
\emph{Encoding.} Let \(L^\sharp\) be the one-sorted vocabulary with unary
relation symbols \(\mathrm{Pt}\), \(\mathrm{Op}\), a binary relation symbol
\(E\), and a unary relation symbol \(Q_p\) for each \(p\in\mathrm{Prop}\).
It is countable. For a two-sorted structure \(S=(X',O',E',(Q'_p)_p)\) let
\(S^\sharp\) have domain \((\{0\}\times X')\cup(\{1\}\times O')\), with
\(\mathrm{Pt}=\{0\}\times X'\), \(\mathrm{Op}=\{1\}\times O'\),
\(E=\{((0,x),(1,u)):(x,u)\in E'\}\) and
\(Q_p=\{(0,x):x\in Q'_p\}\). The tags make the two parts disjoint even when
an open is itself a point. Translate an \(L_2\)-formula \(\varphi\) to
\(\varphi^\sharp\), keeping its variables: \(x=y\), \(U=V\) and truth are
unchanged, \(x\varepsilon U\) becomes \(E(x,U)\), \(P_p(x)\) becomes
\(Q_p(x)\), connectives commute with the translation, and
\[
 (\exists x\,\psi)^\sharp=\exists x\,(\mathrm{Pt}(x)\land\psi^\sharp),\qquad
 (\exists U\,\psi)^\sharp=\exists U\,(\mathrm{Op}(U)\land\psi^\sharp).
\]
Universal quantifiers are read as \(\neg\exists\neg\), so they translate
to the relativized universal forms. For an assignment \(s\) respecting
sorts let \(s^\sharp\) send each point variable \(x\) to \((0,s(x))\) and
each open variable \(U\) to \((1,s(U))\).

\emph{Correspondence.} For every two-sorted \(S\), every \(L_2\)-formula
\(\varphi\) and every such \(s\),
\(S\models\varphi[s]\) iff \(S^\sharp\models\varphi^\sharp[s^\sharp]\).
The proof is by induction on \(\varphi\). For atoms it is the definition
of \(S^\sharp\), using that \(a\mapsto(0,a)\) and \(u\mapsto(1,u)\) are
injective, so equalities are preserved and reflected. The connective steps
are immediate. For \(\exists x\,\psi\): \(S\models\exists x\,\psi[s]\) iff
some \(a\in X'\) has \(S\models\psi[s(x/a)]\), iff (induction) some
\(a\in X'\) has \(S^\sharp\models\psi^\sharp[s^\sharp(x/(0,a))]\). The
elements of \(S^\sharp\) satisfying \(\mathrm{Pt}\) are exactly the pairs
\((0,a)\) with \(a\in X'\), so this holds iff
\(S^\sharp\models\exists x(\mathrm{Pt}(x)\land\psi^\sharp)[s^\sharp]\).
The open quantifier step is the same with \(\mathrm{Op}\) and the tag
\(1\).

\emph{Downward step.} Apply \cref{thm:downward-ls} to \(M^\sharp\) and
the countable set
\[
 A=(\{0\}\times A_{\mathrm{pt}})\cup(\{1\}\times A_{\mathrm{op}}).
\]
It gives \(N\preccurlyeq M^\sharp\) containing \(A\) with
\(\lvert\dom N\rvert\le\max(\lvert A\rvert,\lvert L^\sharp\rvert,\aleph_0)
=\aleph_0\). Put \(X_0=\{x:(0,x)\in\dom N\}\) and
\(O_0=\{u:(1,u)\in\dom N\}\). These are countable and contain
\(A_{\mathrm{pt}}\) and \(A_{\mathrm{op}}\). They are nonempty: \(X\ne
\varnothing\) and \(X\in\tau\), so \(M^\sharp\models\exists v\,\mathrm{Pt}(v)\)
and \(M^\sharp\models\exists v\,\mathrm{Op}(v)\), and elementarity gives
witnesses in \(\dom N\). The relations of \(N\) are the restrictions of
those of \(M^\sharp\), so \(N=(M\restriction(X_0,O_0))^\sharp\).

\emph{Decoding.} For \(\varphi\in L_2\) and parameters
\(\bar a\in X_0\), \(\bar u\in O_0\), with \(s\) the corresponding
assignment,
\[
 M\models\varphi[s]
 \iff M^\sharp\models\varphi^\sharp[s^\sharp]
 \iff N\models\varphi^\sharp[s^\sharp]
 \iff M\restriction(X_0,O_0)\models\varphi[s],
\]
by the correspondence for \(M\), elementarity of \(N\) (all values of
\(s^\sharp\) lie in \(\dom N\)), and the correspondence for
\(M\restriction(X_0,O_0)\).
\end{proof}
\end{annexitem}

\begin{annexitem}{annex:p03-w03-p03-051-extent}{thm:p03-w03-p03-051-extent}{Basis axioms, the trace isomorphism, and basis invariance}
\begin{proof}
Write \(S_0=M\restriction(X_0,O_0)\).

\emph{(1)} The four sentences \(\mathrm{Ext}\), \(\mathrm{Union}\),
\(\mathrm{Empty}\), \(\mathrm{Bas}\) of
\cref{def:p03-topological-basis-axioms} hold in \(M\): two members of
\(\tau\) with the same elements are equal; \(X\in\tau\) contains every
point; \(\varnothing\in\tau\) contains none; and for \(x\in U\cap V\) with
\(U,V\in\tau\), the set \(W=U\cap V\in\tau\) contains \(x\) and lies in
both. By \cref{thm:p03-w03-p03-051-two-sorted-ls} they hold in \(S_0\),
where point quantifiers range over \(X_0\) and open quantifiers over
\(O_0\). \(\mathrm{Ext}\) in \(S_0\) says that \(u,v\in O_0\) are equal when
\(u\cap X_0=v\cap X_0\); so \(e\) is injective. \(\mathrm{Empty}\) gives
\(u\in O_0\) with \(e(u)=\varnothing\), so \(\varnothing\in\sigma_0\).
\(\mathrm{Union}\) gives each \(x\in X_0\) some \(u\in O_0\) with
\(x\in e(u)\), so \(\sigma_0\) covers \(X_0\). \(\mathrm{Bas}\) says that if
\(x\in e(u)\cap e(v)\) with \(x\in X_0\), there is \(w\in O_0\) with
\(x\in w\) and every \(z\in X_0\cap w\) in \(u\cap v\); that is,
\(x\in e(w)\subseteq e(u)\cap e(v)\). These are the three conditions on a
basis in \cref{def:p03-topological-carrier}.

\emph{(2)} Let \(B_0=(X_0,\sigma_0,\nu_0)\) be the basoid model and
\(h\) the pair of maps \(\mathrm{id}_{X_0}\) and \(e:O_0\to\sigma_0\).
Both are bijections. For \(x\in X_0\) and \(u\in O_0\), \(x\in u\) iff
\(x\in u\cap X_0=e(u)\), and \(x\in\nu(p)\) iff \(x\in\nu_0(p)\). So \(h\)
is an isomorphism of two-sorted structures from \(S_0\) onto \(B_0\). An
isomorphism preserves and reflects every \(L_2\)-formula under
transported assignments: atoms by the preceding sentence and bijectivity,
connectives trivially, and a quantifier step because a witness \(c\) on
one side gives the witness \(h(c)\) on the other, and every element of
the target sort is some \(h(c)\). Now let \(\varphi\in L_t\),
\(\bar a\in X_0\), \(\bar u\in O_0\). Since \(L_t\subseteq L_2\),
\[
\begin{aligned}
 M\models\varphi[\bar a,\bar u]
 &\iff S_0\models\varphi[\bar a,\bar u]\\
 &\iff B_0\models\varphi[\bar a,e(\bar u)]\\
 &\iff (X_0,\tau_0,\nu_0)\models\varphi[\bar a,e(\bar u)],
\end{aligned}
\]
by \cref{thm:p03-w03-p03-051-two-sorted-ls}, the isomorphism, and
\cref{thm:p03-topological-basis} applied to \(B_0\) with the free open
parameters given the same subsets \(e(\bar u)\) of \(X_0\) on both sides.
For sentences there are no parameters, so the two topological models are
\(L_t\)-equivalent.

\emph{(3)} Every \(u\in O_0\) is in \(\tau\), so \(e(u)=u\cap X_0\) is
open in the subspace topology; unions of such traces are traces of unions,
so \(\tau_0\subseteq\{W\cap X_0:W\in\tau\}\). Conversely let \(W\in\tau\)
and \(x\in W\cap X_0\). Since \(A_{\mathrm{op}}\) is a basis of \(\tau\),
some \(u\in A_{\mathrm{op}}\subseteq O_0\) has \(x\in u\subseteq W\), and
then \(x\in e(u)\subseteq W\cap X_0\). So \(W\cap X_0\) is the union of the
members of \(\sigma_0\) it contains, hence lies in \(\tau_0\).
\end{proof}
\end{annexitem}

\begin{annexitem}{annex:p03-w03-p03-051-real-line}{thm:p03-w03-p03-051-real-line}{The countable piece of the line}
\begin{proof}
Apply \cref{thm:p03-w03-p03-051-two-sorted-ls} to
\(M=(\mathbb R,\tau_{\mathbb R},\nu_\varnothing)\) with
\(A_{\mathrm{pt}}=\mathbb Q\) and \(A_{\mathrm{op}}=\mathcal I_{\mathbb Q}\).
Both are countable: apart from the empty set,
\(\mathcal I_{\mathbb Q}\) is the image of a subset of
\(\mathbb Q\times\mathbb Q\). This gives countable \(X_0\supseteq\mathbb Q\)
and \(O_0\supseteq\mathcal I_{\mathbb Q}\).

The family \(\mathcal I_{\mathbb Q}\) contains the empty set and is a
basis of \(\tau_{\mathbb R}\): if
\(x\in W\in\tau_{\mathbb R}\), some \(\delta>0\) has
\((x-\delta,x+\delta)\subseteq W\), and density of \(\mathbb Q\)
gives rationals \(a\in(x-\delta,x)\) and \(b\in(x,x+\delta)\), so
\(x\in(a,b)\subseteq W\). By \cref{thm:p03-w03-p03-051-extent}(3) the
topology \(\tau_0\) generated by the traces is
\(\tau_{\mathbb R}\restriction X_0\), and \(\nu_0=\nu_\varnothing\). Part~(2)
of the same lemma, with \(\bar u\) from
\(\mathcal I_{\mathbb Q}\subseteq O_0\) and \(e(\bar u)=\bar u\cap X_0\),
is assertion~(1).

\emph{\(X_0\) is disconnected.} \(X_0\) is countable and \(\mathbb R\) is
not (Cantor), so there is \(r\in\mathbb R\setminus X_0\). Density gives
rationals \(q_1<r<q_2\); both lie in \(X_0\). Put
\(L=(-\infty,r)\cap X_0\) and \(R=(r,\infty)\cap X_0\). They are open in
the subspace topology, disjoint, and \(L\cup R=X_0\) because \(r\notin X_0\).
So \(L\) is open and closed, contains \(q_1\), and misses \(q_2\); it is
neither \(\varnothing\) nor \(X_0\).

\emph{\(\mathbb R\) is connected.} Suppose \(\mathbb R=A\cup B\) with \(A,B\)
open, disjoint and nonempty; choose \(a\in A\), \(b\in B\), and by symmetry
assume \(a<b\). Let \(T=\{t\in[a,b]:[a,t]\subseteq A\}\), which contains
\(a\) and is bounded by \(b\), and let \(s=\sup T\), so \(a\le s\le b\).
Every \(t\in[a,s)\) lies below some member of \(T\), hence
\([a,s)\subseteq A\). If \(s\in A\), choose \(\delta>0\) with
\((s-\delta,s+\delta)\subseteq A\). Then \(b\ge s+\delta\),
since otherwise \(b\in(s-\delta,s+\delta)\subseteq A\). Hence
\([a,s+\delta/2]\subseteq A\) and \(s+\delta/2\in T\), contrary
to \(s=\sup T\). If \(s\in B\), then \(s\ne a\), so \(s>a\); choose
\(\delta>0\) with \((s-\delta,s+\delta)\subseteq B\) and a
point \(t\) with \(\max(a,s-\delta)<t<s\). Then \(t\in A\cap B\), a
contradiction. So no such splitting exists, and the only subsets of
\(\mathbb R\) that are open with open complement are \(\varnothing\) and
\(\mathbb R\).
\end{proof}
\end{annexitem}

\begin{annexitem}{annex:p03-w03-p03-051-connectedness}{thm:p03-w03-p03-051-connectedness}{From one equivalent pair to non-axiomatizability}
\begin{proof}
Let \(X_0\) be as in \cref{thm:p03-w03-p03-051-real-line}, and put
\(M_1=(\mathbb R,\tau_{\mathbb R},\nu_\varnothing)\) and
\(M_2=(X_0,\tau_{\mathbb R}\restriction X_0,\nu_\varnothing)\). Both belong
to \(K_{\mathbb R}\) (\(X_0\supseteq\mathbb Q\) is nonempty), and by
part~(1) of that theorem with no parameters they satisfy the same
\(L_t\)-sentences.

(1) Suppose \(\Gamma\) axiomatizes connectedness relative to
\(K_{\mathbb R}\). The space of \(M_1\) is connected, so
\(M_1\models\Gamma\); hence \(M_2\models\Gamma\), and \(M_2\) would be
connected, contrary to part~(2) of the theorem. If \(\Gamma\) axiomatizes
the failure of connectedness, then \(M_2\models\Gamma\), hence
\(M_1\models\Gamma\), and \(M_1\) would be disconnected, again a
contradiction.

(2) A sentence \(\chi\) expressing connectedness relative to
\(K_{\mathbb R}\) would make \(\{\chi\}\) an axiomatizing set, excluded by
(1). A sentence expressing it relative to all topological models would in
particular express it relative to \(K_{\mathbb R}\).

(3) If every sentence of \(\Delta\) is true in \(M_1\), it is true in
\(M_2\) by the \(L_t\)-equivalence, so both models lie in the subclass of
\(K_{\mathbb R}\) defined by \(\Delta\). The arguments for (1) and (2) used
only these two models, so they apply verbatim to that subclass.
\end{proof}
\end{annexitem}

\begin{annexitem}{annex:p03-w03-p03-051-compactness}{thm:p03-w03-p03-051-compactness}{The unit interval}
\begin{proof}
Let \(I=[0,1]\) with \(\tau_I=\tau_{\mathbb R}\restriction I\), and put
\(M=(I,\tau_I,\nu_\varnothing)\). The traces
\(\mathcal I'=\{(a,b)\cap I:(a,b)\in\mathcal I_{\mathbb Q}\}\) are a
countable basis of \(\tau_I\): an open set of \(\tau_I\) is \(W\cap I\)
with \(W\in\tau_{\mathbb R}\), and for \(x\in W\cap I\) the basis property
of \(\mathcal I_{\mathbb Q}\), proved in the preceding item, gives
\(x\in(a,b)\subseteq W\), so \(x\in(a,b)\cap I\subseteq W\cap I\). Apply
\cref{thm:p03-w03-p03-051-two-sorted-ls} with
\(A_{\mathrm{pt}}=\mathbb Q\cap I\) and \(A_{\mathrm{op}}=\mathcal I'\), and
then \cref{thm:p03-w03-p03-051-extent}(2),(3). They give a countable
\(Z_0\) with \(\mathbb Q\cap I\subseteq Z_0\subseteq I\) such that \(M\) and
\((Z_0,\tau_I\restriction Z_0,\nu_\varnothing)\) are \(L_t\)-equivalent;
and \(\tau_I\restriction Z_0=\tau_{\mathbb R}\restriction Z_0\).

\emph{\(I\) is compact.} Let \(\mathcal C\subseteq\tau_I\) cover \(I\), and
let \(T\) be the set of \(t\in I\) such that finitely many members of
\(\mathcal C\) cover \([0,t]\). Then \(0\in T\); let \(s=\sup T\in I\).
Some \(C\in\mathcal C\) contains \(s\), and since \(C\) is open in \(I\)
there is \(\delta>0\) with
\((s-\delta,s+\delta)\cap I\subseteq C\). Since \(s=\sup T\), some
\(t\in T\) has \(s-\delta<t\le s\). Finitely many
members cover \([0,t]\); adding \(C\) covers
\([0,t']\) with \(t'=\min(s+\delta/2,1)\), since
\([t,t']\subseteq(s-\delta,s+\delta)\cap I\). So \(t'\in T\) and
\(t'\le s\), which forces \(t'=1\). Hence \(1\in T\), a finite subcover.

\emph{\(Z_0\) is not compact.} \(Z_0\) is countable and \(I\) is not, so
some \(r\in I\setminus Z_0\); \(r\) is irrational, so \(0<r<1\). For
\(n\ge1\) let \(C_n=Z_0\setminus[r-1/n,r+1/n]\), open in the subspace
topology of \(Z_0\). They increase with \(n\), and they cover \(Z_0\): every
\(z\in Z_0\) differs from \(r\), and \(z\in C_n\) once \(1/n<\lvert z-r\rvert\).
A finite subfamily has union \(C_m\) for its largest index \(m\). Density
gives a rational \(q\in I\) with \(0<\lvert q-r\rvert<1/m\); it lies in
\(Z_0\setminus C_m\). So no finite subfamily covers \(Z_0\).

The two models are in \(K_{\mathbb R}\) and are \(L_t\)-equivalent; the
argument of \cref{thm:p03-w03-p03-051-connectedness}(1),(2), with
compactness in place of connectedness, gives the remaining assertions.
\end{proof}
\end{annexitem}

\begin{annexitem}{annex:p03-w03-p03-051-conn-basis}{thm:p03-w03-p03-051-conn-basis}{The connectedness sentence on the basis and on the topology}
\begin{proof}
By \cref{thm:p03-topological-property-formulas}, \(\mathrm{Conn}\) holds in a
topological model exactly when its space is connected; so it holds in
\((\mathbb R,\tau_{\mathbb R},\nu_\varnothing)\) by
\cref{thm:p03-w03-p03-051-real-line}(2). It is an \(L_2\)-sentence, so
\cref{thm:p03-w03-p03-051-two-sorted-ls} gives it in
\(S_0=(\mathbb R,\tau_{\mathbb R},\nu_\varnothing)\restriction(X_0,O_0)\),
and the isomorphism \(h\) of the proof of
\cref{thm:p03-w03-p03-051-extent}(2) carries it to the basoid model
\((X_0,\sigma_0,\nu_\varnothing)\). In the generated model
\((X_0,\tau_{\mathbb R}\restriction X_0,\nu_\varnothing)\) the space is
disconnected (\cref{thm:p03-w03-p03-051-real-line}(2)), so the same
characterization makes \(\mathrm{Conn}\) false there. Concretely, in the
basoid model the open quantifiers range over \(\sigma_0\), and the
sentence says that whenever \(e(u)\) and \(e(v)\) are complementary in
\(X_0\), one of them is \(\varnothing\) or \(X_0\).
\end{proof}
\end{annexitem}

\begin{annexitem}{annex:p03-w03-p03-051-modal}{thm:p03-w03-p03-051-modal}{Pointwise modal agreement}
\begin{proof}
Let \(a\in X_0\). By \cref{thm:p03-topological-translation-homeomorphism},
\((\mathbb R,\tau_{\mathbb R},\nu_\varnothing),a\models\varphi\) iff
\((\mathbb R,\tau_{\mathbb R},\nu_\varnothing)\models\mathrm{ST}_x(\varphi)[a]\),
and likewise for the countable model. The translation is an
\(L_t\)-formula with the single free point variable \(x\)
(\cref{def:p03-topological-modal-translation}), so
\cref{thm:p03-w03-p03-051-real-line}(1) with \(\bar a=a\) and no open
parameters gives the displayed equivalence.

For the second assertion, \(\varphi\) is true at every point of a model iff
the model satisfies the \(L_t\)-sentence \(\forall x\,\mathrm{ST}_x(\varphi)\),
by the same translation theorem. The two models of
\cref{thm:p03-w03-p03-051-real-line} are \(L_t\)-equivalent, so \(\varphi\)
is true at every point of the first iff at every point of the second. The
first space is connected and the second is not, so \(\varphi\) does not
separate them.
\end{proof}
\end{annexitem}
